\section{Enfoque supervisado con CNN ResNet18}
\label{chap:cnn_enfoque}

La CNN constituye la referencia supervisada del trabajo. Su función no es competir con todas las arquitecturas posibles, sino representar un enfoque estándar para aprender a partir de imágenes etiquetadas. En este caso, los ejemplos disponibles se usan para actualizar pesos y producir un modelo específico para la tarea.

\subsection{Arquitectura y entrada}

El modelo ejecutado es una ResNet18 preentrenada \cite{he2016}. La cabeza final se sustituye por tres salidas, una para cada hallazgo morfológico: RBBB, elevación ST e inversión de T. La entrada se procesa como imagen RGB de 224 por 224 píxeles. En el simulador, la etiqueta positiva de paciente se obtiene cuando los tres hallazgos están presentes.

Esta formulación condiciona toda la evaluación real. La CNN no aprende directamente la etiqueta clínica \texttt{clinical\_brugada}. Aprende a detectar tres rasgos morfológicos y después deriva una decisión positiva cuando los tres aparecen simultáneamente. Por ello, los datos reales del HUCA se usan como conjunto de prueba clínica, no como conjunto supervisado de entrenamiento morfológico.

\begin{figure}[H]
  \centering
  \safeincludegraphics[width=0.96\textwidth]{results/cnn_convolutional_pipeline.drawio.png}
  \caption{Uso supervisado de los ejemplos etiquetados en la CNN. Las muestras de entrenamiento actualizan los pesos de ResNet18 y la muestra reservada produce una predicción binaria junto con su mapa Grad CAM.}
  \label{fig:cnn_convolutional_pipeline}
\end{figure}

\subsection{Entrenamiento}

La pérdida es entropía cruzada binaria multietiqueta. Cuando el pliegue lo requiere, la clase positiva se pondera automáticamente a partir de los ejemplos de entrenamiento. El optimizador es AdamW \cite{loshchilov2019}, con tasa de aprendizaje \(3\times10^{-4}\), lote de 32 imágenes y 20 épocas máximas.

\begin{table}[H]
  \centering
  \small
  \begin{tabularx}{0.92\textwidth}{p{4cm}X}
    \toprule
    Propiedad & Valor \\
    \midrule
    Arquitectura & ResNet18 con cabeza multietiqueta de tres salidas \\
    Inicialización & Pesos preentrenados de ResNet18 \\
    Entrada & RGB, 224 por 224 píxeles \\
    Pérdida & Entropía cruzada binaria multietiqueta con ponderación automática \\
    Optimizador & AdamW \\
    Tasa de aprendizaje & \(3\times10^{-4}\) \\
    Épocas máximas & 20 \\
    Tamaño de lote & 32 \\
    Valores de \(K\) & \(K=\{2,4,8,16,32\}\) \\
    Semillas & 42, 123 y 2026 \\
    Umbral & Seleccionado por validación con exactitud equilibrada por paciente \\
    Agregación & Sin agregación multiderivación; se usa únicamente V1 antes del umbral \\
    \bottomrule
  \end{tabularx}
  \caption{Configuración CNN utilizada en los experimentos finales.}
  \label{tab:modelo_cnn_final}
\end{table}

\subsection{Uso de datos reales en la CNN}

Los datos reales del HUCA no contienen anotaciones independientes de RBBB, elevación ST e inversión de T para cada imagen. Solo incorporan la referencia clínica de caso Brugada confirmado frente a normal. Esto impide entrenar, con la arquitectura definida, una CNN supervisada con esa cohorte real y evaluarla después sobre la misma fuente manteniendo la tarea morfológica.

La evaluación real de la CNN se plantea por tanto como transferencia desde el simulador. En cada pliegue, los \(K\) pacientes que actualizan los pesos proceden del conjunto sintético QRS/ST. La muestra evaluada pertenece a los datos reales del HUCA y su resultado se compara con la etiqueta clínica disponible. Esta decisión no inventa etiquetas morfológicas en el dominio real, pero hace que la lectura sea necesariamente una prueba de transferencia sintético-real.

Entrenar una CNN supervisada directamente con los datos reales del HUCA sería una línea distinta. Requeriría disponer de anotaciones morfológicas clínicas por imagen o rediseñar el modelo como clasificador binario directo de \texttt{clinical\_brugada}. Esa segunda opción mediría otra tarea, más próxima a predicción clínica agregada que a detección de rasgos visuales.

\subsection{Validación por paciente}

El entrenamiento y la evaluación se organizan mediante validación leave-one-out. En cada pliegue se reserva un paciente completo y el modelo se ajusta con un subconjunto de \(K\) pacientes del resto. Esto evita que un mismo paciente aparezca simultáneamente en entrenamiento y prueba. En los datos reales del HUCA esta restricción es indispensable porque la etiqueta clínica pertenece al paciente, no a una imagen aislada.

La evaluación registra las probabilidades de la imagen V1, la predicción por paciente, el umbral elegido, la matriz de confusión y las métricas finales. También conserva paneles Grad CAM para inspección cualitativa.

\subsection{Selección de umbral}

La ResNet18 se entrena con tres salidas, una para RBBB, otra para elevación ST y otra para inversión de T. Cada salida produce un logit, que se transforma mediante una sigmoide en una probabilidad. Por tanto, el umbral no se aplica sobre una clase Brugada directa, sino sobre cada hallazgo morfológico. En esta versión experimental solo se usa V1, de modo que las probabilidades de V1 se umbralizan directamente por hallazgo.

La decisión de paciente se obtiene después de aplicar el umbral a V1. Un hallazgo se considera presente si su probabilidad alcanza el umbral elegido. La predicción positiva de Brugada se deriva con la misma regla utilizada durante el experimento: RBBB, elevación ST e inversión de T deben estar presentes simultáneamente.

El valor del umbral se escoge dentro de cada pliegue usando solo pacientes de validación, nunca el paciente reservado para prueba. La búsqueda prueba umbrales candidatos y conserva el que maximiza la exactitud equilibrada de la etiqueta derivada a nivel de paciente. Si la validación del pliegue no contiene ambas clases, se mantiene 0,5. El paciente de prueba se evalúa siempre con un criterio fijado antes de consultar su etiqueta.

\subsection{Adaptación de dominio}

Además de la transferencia directa, se probaron variantes de adaptación de dominio no supervisada. En todas ellas el origen etiquetado sigue siendo el simulador QRS/ST y los datos reales del HUCA se usan como dominio objetivo sin etiquetas morfológicas. El objetivo no es convertir la cohorte real en un conjunto de entrenamiento clínico, sino reducir la discrepancia entre las representaciones aprendidas en imágenes sintéticas y las observadas en ECG reales.

El preentrenamiento autosupervisado usa SimCLR sobre imágenes sin etiquetas de los datos reales del HUCA \cite{chen2020simclr}. El encoder recibe dos vistas aumentadas de una misma imagen y aprende a acercar sus representaciones frente a otras imágenes del lote. Después, esos pesos inicializan la CNN supervisada sobre el simulador.

CORAL añade una pérdida que aproxima las covarianzas de las características fuente y destino, de modo que las activaciones de ambos dominios tengan estadísticas de segundo orden más parecidas \cite{sun2016deepcoral}. MMD mide la discrepancia entre distribuciones de características mediante kernels RBF y penaliza que las medias embebidas de origen y destino permanezcan separadas \cite{long2015dan}. DANN incorpora un clasificador de dominio y una capa de inversión de gradiente, de forma que el extractor aprende características útiles para la tarea supervisada pero poco informativas sobre si una imagen procede del simulador o de HUCA \cite{ganin2016}.

\begin{table}[H]
  \centering
  \small
  \begin{tabularx}{0.94\textwidth}{p{2.6cm}X}
    \toprule
    Método & Idea principal \\
    \midrule
    SSL & Preentrenar el encoder con imágenes reales sin etiquetas mediante aprendizaje contrastivo. \\
    CORAL & Alinear covarianzas de características entre simulador y HUCA. \\
    MMD & Reducir una distancia de distribución con kernels RBF. \\
    DANN & Aprender características invariantes al dominio con inversión de gradiente. \\
    \bottomrule
  \end{tabularx}
  \caption{Métodos de adaptación de dominio evaluados para la CNN.}
  \label{tab:cnn_domain_methods}
\end{table}

\subsection{Grad CAM}

Grad CAM se utiliza como diagnóstico cualitativo, no como explicación clínica suficiente \cite{selvaraju2017}. El mapa indica regiones que influyen en el logit del modelo, pero no prueba que la red haya aprendido un criterio médico correcto. En esta memoria se usa para comprobar si las activaciones se concentran en zonas plausibles del trazado y para detectar errores en los que una atención visual razonable termina en una decisión equivocada.

\begin{figure}[H]
  \centering
  \safeincludegraphics[width=0.92\textwidth]{results/cnn_simulator_qrs_gradcam_example.png}
  \caption{Ejemplo Grad CAM del modelo CNN sobre el conjunto sintético.}
  \label{fig:gradcam_sintetico_modelos}
\end{figure}

\subsection{Qué se mide}

La evaluación posterior mide tres cuestiones. La primera es si una CNN entrenada con pocos ejemplos aprende una regla visual en el simulador QRS/ST. La segunda es si ese modelo transferido desde el simulador mantiene utilidad cuando la etiqueta pasa a ser clínica y los ECG proceden de pacientes reales. La tercera es si la adaptación de dominio no supervisada reduce de forma apreciable esa brecha. Estas preguntas se analizan después en los dominios sintético y real.
